\chapter{Vector calculus and the scalar product}\label{C03}
\begin{repere}{Chapter path}
\textbf{Objectives.} Choose between projection, coordinates and angle; calculate a length; characterise orthogonality and a locus.

\textbf{Prerequisites.} Vectors, Chasles's relation, Pythagoras, cosine in a right triangle, orthonormal frame.

\textbf{Yearly reference point.} 6--15 October: 8 hours\end{repere}
\section*{Diagnostic: six questions, 10 minutes}\label{C03-D}
Answer without consulting the lesson. Mark the questions for which help is needed. This diagnostic is not graded.
\begin{enumerate}
\item For $A(1;2)$ and $B(4;-2)$, calculate $\vect{AB}$.
\item Calculate the norm of $(3;-4)$.
\item Simplify $\vect{AB}+\vect{BC}$.
\item Expand $(a+b)^2$.
\item In a right triangle, an angle has adjacent side 3 and hypotenuse 5. Give its cosine.
\item What is the midpoint of $A(-2;1)$ and $B(4;3)$?
\end{enumerate}
\paragraph{After the diagnostic.} Revisit the flagged topics with the teacher; redo a similar question, then explain the rule used.
\paragraph{Speaking and logic.} Compare, on the same problem, coordinates and vector calculus without a frame. Justify why your method uses the data most directly.
\clearpage
\section*{Investigate before formalising}\label{C03-RECH}
Construct three points M such that $\vect{MA}\cdot\vect{MB}=0$ for a segment AB of 6 cm. Conjecture the locus, then produce a proof that includes A and B.
\paragraph{Organisation.} Individual investigation 5 min, comparison in pairs 5 min, class discussion 5 min. Keep a record of one attempt, even an unfinished one. A conjecture must then be justified.
\section{Projection and angle}\label{C03-U1}
\begin{repere}{Definition and results}
For two nonzero vectors $\vec u,\vec v$, their scalar product is the real number $\vec u\cdot\vec v=\norm{\vec u}\norm{\vec v}\cos\theta$, where $\theta\in[0;\pi]$ is their unoriented angle. With a zero vector, the product is 0. If $H$ is the projection of B onto $(OA)$, $\vect{OA}\cdot\vect{OB}$ equals $OA\times OH$ when H is on the ray $[OA)$, and its opposite otherwise.
\end{repere}
\paragraph{Status and justification.} Definition; the equivalence with the projection follows from the cosine in right triangles, with a negative algebraic projection for an obtuse angle.
\exemple{For two vectors of norms 3 and 4 forming $60^\circ$, the product is $\dfrac{3\times 4\times 1}{2}=6$. With $120^\circ$, it is $-6$. If $A(3;0)$ and $B(-2;4)$, the projection is $H(-2;0)$ and the product is $-6$.}\label{C03-EX1}
\par\medskip\noindent\begin{minipage}{\linewidth}\centering
\begin{tikzpicture}[x=1cm,y=1cm,>=Stealth]\draw[->](-3,0)--(4,0)node[right]{$x$};\draw[->](0,-.6)--(0,4.8)node[above]{$y$};\draw[->,very thick,bleu](0,0)--(3,0);\draw[->,very thick,vert](0,0)--(-2,4);\draw[dashed](-2,4)--(-2,0);\draw(-2,.25)--(-1.75,.25)--(-1.75,0);\node[below right]at(0,0){$O$};\node[below]at(3,0){$A(3;0)$};\node[below]at(-2,0){$H(-2;0)$};\node[above left]at(-2,4){$B(-2;4)$};\end{tikzpicture}
\par\medskip The projection H lies on the negative side of the horizontal axis. Hence $\vect{OA}\cdot\vect{OB}=3\times(-2)=-6$.\end{minipage}\par\medskip

\begin{erreur}The result is a number, not a vector. A right angle corresponds to 0, not to 1.\end{erreur}
\application{Two vectors have norms 2 and 5 and form $90^\circ$, then $180^\circ$. Calculate the products.}\label{C03-AP1}
\section{Coordinates, norm and orthogonality}\label{C03-U2}
\begin{repere}{Definition and results}
In an orthonormal basis, for $\vec u=(x;y)$ and $\vec v=(x\prime;y\prime)$, $\vec u\cdot\vec v=xx\prime+yy\prime$ and $\norm{\vec u}^2=x^2+y^2$. The vectors are orthogonal if and only if their product is zero. For an orthonormal basis $(\vec i,\vec j)$, $x=\vec u\cdot\vec i$ and $y=\vec u\cdot\vec j$.
\end{repere}
\paragraph{Status and justification.} Properties proved using bilinearity: expand $(x\vec i+y\vec j)\cdot(x\prime\vec i+y\prime\vec j)$ and use $\vec i\cdot\vec i=\vec j\cdot\vec j=1$, $\vec i\cdot\vec j=0$.
\exemple{$\vec u=(2;3)$, $\vec v=(3;-2)$: product $6-6=0$. The two vectors are orthogonal. Their norm is $\sqrt{13}$. The product of $(2;3)$ with $(1;0)$ is 2: we recover the x-coordinate.}\label{C03-EX2}
\begin{erreur}The usual formula assumes an orthonormal basis. Collinearity and orthogonality are different properties.\end{erreur}
\application{For which real number t are the vectors $(t;2)$ and $(3;-6)$ orthogonal?}\label{C03-AP2}
\section{Vector calculus and identities}\label{C03-U3}
\begin{repere}{Definition and results}
The scalar product is symmetric and bilinear: $\vec u\cdot(\vec v+\vec w)=\vec u\cdot\vec v+\vec u\cdot\vec w$ and $(\lambda\vec u)\cdot\vec v=\lambda(\vec u\cdot\vec v)$. Hence $\norm{\vec u\pm\vec v}^2=\norm{\vec u}^2\pm2\vec u\cdot\vec v+\norm{\vec v}^2$. It follows that $\vec u\cdot \vec v=\dfrac{\norm{\vec u+\vec v}^2-\norm{\vec u}^2-\norm{\vec v}^2}{2}$.
\end{repere}
\paragraph{Status and justification.} Bilinearity accepted from the geometric definition; symmetry comes from the cosine. The identities are proved by expanding $(\vec u\pm\vec v)\cdot(\vec u\pm\vec v)$.
\exemple{If $\norm{\vec u}=3$, $\norm{\vec v}=4$ and $\norm{\vec u+\vec v}=5$, the product is $\dfrac{25-9-16}{2}=0$. We deduce orthogonality without coordinates.}\label{C03-EX3}
\begin{erreur}In general, $\norm{\vec u+\vec v}\ne\norm{\vec u}+\norm{\vec v}$. The identities concern the squares of the norms.\end{erreur}
\application{Given $\norm{\vec u}=2$, $\norm{\vec v}=3$, $\vec u\cdot\vec v=1$. Calculate $\norm{\vec u-\vec v}$.}\label{C03-AP3}
\section{Al-Kashi and choosing a method}\label{C03-U4}
\begin{repere}{Definition and results}
In a triangle ABC, $BC^2=AB^2+AC^2-2AB\,AC\cos\widehat{BAC}$. This formula allows us to calculate the third side from two sides and the included angle, or to find the cosine of an angle from the three sides.
\end{repere}
\paragraph{Status and justification.} By Chasles, $\vect{BC}=\vect{AC}-\vect{AB}$. Expanding its scalar square gives $BC^2=AC^2+AB^2-2\vect{AB}\cdot\vect{AC}$; then introduce the expression with the cosine.
\exemple{$AB=5$, $AC=3$ and $\widehat{BAC}=60^\circ$: $BC^2=25+9-15=19$, so $BC=\sqrt{19}$. With three sides, isolate the cosine before approximating the angle on a calculator, if one is allowed.}\label{C03-EX4}
\begin{erreur}The angle is included between the two sides of the product term. The minus sign is essential.\end{erreur}
\application{A triangle has $AB=AC=4$ and $BC=4\sqrt2$. Determine the angle A.}\label{C03-AP4}
\section{The circle with diameter AB}\label{C03-U5}
\begin{repere}{Definition and results}
For two distinct points A and B with midpoint I, $\vect {MA}\cdot \vect {MB}=MI^2-\dfrac{AB^2}{4}$. The set of points M such that $\vect{MA}\cdot\vect{MB}=0$ is the circle with centre I and radius $\dfrac{AB}{2}$, including the points A and B.
\end{repere}
\paragraph{Status and justification.} Write $\vect{MA}=\vect{MI}+\vect{IA}$ and $\vect{MB}=\vect{MI}-\vect{IA}$. The cross terms cancel; what remains is $MI^2-IA^2$. Since $IA=\dfrac{AB}{2}$, equality to zero is equivalent to the metric equation of the circle.
\exemple{$A(-2;0)$ and $B(2;0)$: I is the origin. The locus is $x^2+y^2=4$. For $M(0;1)$, the product is $1-4=-3$; for $M(0;2)$, it is 0.}\label{C03-EX5}
\par\medskip\noindent\begin{minipage}{\linewidth}\centering
\begin{tikzpicture}[scale=.8]\draw[bleu,thick](0,0)circle(2);\draw(-2,0)--(2,0);\draw[vert,thick](-2,0)--(0,2)--(2,0);\fill(-2,0)circle(2pt)node[left]{$A$};\fill(2,0)circle(2pt)node[right]{$B$};\fill(0,0)circle(2pt)node[below]{$I$};\fill(0,2)circle(2pt)node[above]{$M$};\end{tikzpicture}
\par\medskip Circle with diameter AB: the vectors $\vect{MA}$ and $\vect{MB}$ are orthogonal.\end{minipage}\par\medskip

\begin{erreur}The endpoints belong to the locus even though ``right triangle'' does not apply to a degenerate triangle.\end{erreur}
\application{For $A(0;0)$, $B(6;0)$, determine the points M such that $\vect{MA}\cdot\vect{MB}=7$.}\label{C03-AP5}
\section*{Exercise path}\label{C03-PARC}
Start with exercises 1 and 2. Then choose one of exercises 3 or 4 and do it without help. Exercises 5 to 7 are a reserve for transfer, investigation or programming: the teacher designates one. Exercise 8 is revisited after a delay.
\begingroup
\exercice{C03-E1}{Direct calculations}
In an orthonormal basis of the plane, consider the vectors $\vec u=(1;-2)$ and $\vec v=(4;3)$.
\begin{enumerate}[label=\arabic*.]
\item Calculate the scalar product $\vec u\cdot\vec v$, then the norms $\norm{\vec u}$ and $\norm{\vec v}$.
\item Calculate the coordinates of $\vec u+\vec v$ and deduce $\norm{\vec u+\vec v}^2$.
\item Recover this last result using the identity involving $\vec u\cdot\vec v$.
\end{enumerate}
\exercice{C03-E2}{A right triangle?}
The plane is equipped with an orthonormal frame. Given the points $A(1;1)$, $B(4;2)$ and $C(0;4)$.
\begin{enumerate}[label=\arabic*.]
\item Calculate the coordinates of $\vect{AB}$ and $\vect{AC}$, then their scalar product.
\item Deduce that the triangle $ABC$ is right-angled; state at which vertex.
\item Calculate the lengths of the two perpendicular sides, then the area of the triangle in area units.
\end{enumerate}
\exercice{C03-E3}{Length without a frame}
The quadrilateral $ABCD$ is a parallelogram with $AB=4$, $AD=3$ and $\widehat{BAD}=60^\circ$. The lengths are expressed in the same unit.
\begin{enumerate}[label=\arabic*.]
\item Draw a figure and calculate the scalar product $\vect{AB}\cdot\vect{AD}$.
\item Determine the exact lengths of the diagonals $[AC]$ and $[BD]$. Explain how the vectors corresponding to the diagonals are expressed in terms of $\vect{AB}$ and $\vect{AD}$.
\item State which diagonal is longer and check that this result is consistent with your figure.
\end{enumerate}
\exercice{C03-E4}{Recovering an angle}
Two nonzero vectors $\vec u$ and $\vec v$ satisfy $\norm{\vec u}=2$, $\norm{\vec v}=6$ and $\vec u\cdot\vec v=-6$. Let $\theta$ be their unoriented angle, between $0^\circ$ and $180^\circ$.
\begin{enumerate}[label=\arabic*.]
\item Determine the exact value of $\cos\theta$, then the measure of $\theta$.
\item Explain how the sign of the scalar product made it possible to predict that the angle would be obtuse.
\end{enumerate}
\exercice{C03-E5}{A variable locus}
The points $A$ and $B$ are distinct and fixed. Let $I$ be the midpoint of $[AB]$. For a fixed real number $k$, we look for the set of points $M$ of the plane that satisfy $\vect{MA}\cdot\vect{MB}=k$.
\begin{enumerate}[label=\arabic*.]
\item Express this scalar product in terms of $MI$ and $AB$ using Chasles's relation and a scalar identity.
\item Describe precisely the set sought according to the value of $k$. For each case, state whether it is empty, a single point or a circle; then give its characteristics.
\item Explain the role of the fact that a squared distance is positive or zero.
\end{enumerate}
\exercice{C03-E6}{Work done by a force}
An object moves in a straight line from $A$ to $B$ over a distance of 4 m. It is subject to a constant force $\vec F$ of magnitude 50 N. The work done by this force is defined by $W=\vec F\cdot\vect{AB}$ and is expressed in joules.
\begin{enumerate}[label=\arabic*.]
\item Calculate $W$ when the angle between $\vec F$ and $\vect{AB}$ is $60^\circ$.
\item Repeat the calculation for an angle of $120^\circ$.
\item Compare the signs of the two works and explain, in each situation, whether the force helps the displacement or opposes it.
\end{enumerate}
\exercice{C03-E7}{Contrapositive}
We wish to prove the following statement, valid for two vectors of the plane:\[\norm{\vec u+\vec v}^2\ne\norm{\vec u}^2+\norm{\vec v}^2\quad\Longrightarrow\quad \vec u\text{ and }\vec v\text{ are not orthogonal}.\]
\begin{enumerate}[label=\arabic*.]
\item State the contrapositive of this implication.
\item Prove this contrapositive using a scalar identity, then conclude about the initial statement.
\end{enumerate}
\exercice{C03-E8}{Delayed review}
A triangle $ABC$ satisfies $AB=3$, $AC=4$ and $BC=5$. We look for $\vect{AB}\cdot\vect{AC}$ without using coordinates.
\begin{enumerate}[label=\arabic*.]
\item Calculate this scalar product using a formula involving the three lengths.
\item Recover the result by determining the nature of the triangle. Name the theorem used and its numerical hypothesis.
\end{enumerate}
\paragraph{Homework.} Do only the task designated at the end of the session, in the dated progression: 10--20 min at most. The exercises depend on the topics already studied, so some questions are deferred. Revisit exercise 8 one week after the chapter (10 min). If the allotted time is exceeded, note the difficulty and bring your attempts; do not pile up the whole reserve.
\endgroup
\Needspace{12\baselineskip}\section*{To do in the CQFD textbook}
These exercises are done in your textbook. The progression indicates the session; they replace the homework planned for that day (10 to 20 minutes in total).
\paragraph{CQFD, exercise 9 p.~250, questions a and b.}\label{C03-M1} Calculate the scalar product in the orthonormal basis and detail the two products of coordinates.
\paragraph{CQFD, exercise 11 p.~250, questions 1, 2 and 3.}\label{C03-M2} Use in turn the coordinates, Chasles's relation and the fact that I is the midpoint.
\begingroup\small
\clearpage\section*{MCQ 1: understand and choose}\label{C03-Q1}
Five single-answer questions for practice. Choose one answer per question; then explain one distractor. This chapter format does not reproduce the official EAM format.
\begin{enumerate}
\item $(2;1)\cdot(3;-4)$ equals
\begin{enumerate}[label=\Alph*.,itemsep=2pt]
\item $2$
\item $10$
\item $(6;-4)$
\item $-2$
\end{enumerate}
\item The norm of $(-3;4)$ equals
\begin{enumerate}[label=\Alph*.,itemsep=2pt]
\item $7$
\item $25$
\item $5$
\item $1$
\end{enumerate}
\item Two vectors with zero product are
\begin{enumerate}[label=\Alph*.,itemsep=2pt]
\item equal
\item collinear
\item opposite
\item orthogonal
\end{enumerate}
\item Norms 2 and 3, angle $60^\circ$: product
\begin{enumerate}[label=\Alph*.,itemsep=2pt]
\item $6$
\item $3$
\item $\sqrt3$
\item $0$
\end{enumerate}
\item $\norm{u-v}^2$ equals
\begin{enumerate}[label=\Alph*.,itemsep=2pt]
\item $\norm u^2-\norm v^2$
\item $\norm u^2+\norm v^2$
\item $\norm u^2-2u\cdot v+\norm v^2$
\item $\norm u^2+2u\cdot v+\norm v^2$
\end{enumerate}
\end{enumerate}
\endgroup
\begingroup\small
\Needspace{18\baselineskip}\section*{MCQ 2: understand and choose}\label{C03-Q2}
Five single-answer questions for practice. Choose one answer per question; then explain one distractor. This chapter format does not reproduce the official EAM format.
\begin{enumerate}
\item An obtuse angle between nonzero vectors gives a product that is
\begin{enumerate}[label=\Alph*.,itemsep=2pt]
\item positive
\item zero
\item equal to 1
\item negative
\end{enumerate}
\item In an orthonormal basis, $u\cdot i$ is
\begin{enumerate}[label=\Alph*.,itemsep=2pt]
\item the x-coordinate of u
\item the norm of u
\item the y-coordinate of u
\item always zero
\end{enumerate}
\item The locus $\vect{MA}\cdot\vect{MB}=0$ for $A\ne B$ is
\begin{enumerate}[label=\Alph*.,itemsep=2pt]
\item the perpendicular bisector
\item the line AB
\item the circle with diameter AB
\item the midpoint alone
\end{enumerate}
\item If $u\cdot v=3$, $(2u)\cdot(-v)$ equals
\begin{enumerate}[label=\Alph*.,itemsep=2pt]
\item $-3$
\item $-6$
\item $6$
\item $9$
\end{enumerate}
\item For $AB=AC=2$ and a right angle at A, $BC$ equals
\begin{enumerate}[label=\Alph*.,itemsep=2pt]
\item $4$
\item $2$
\item $8$
\item $2\sqrt2$
\end{enumerate}
\end{enumerate}
\subsection*{True or false? Justify.}\label{C03-VF}
\begin{enumerate}
\item The scalar product of two vectors is always positive.
\item The zero vector is orthogonal to every vector.
\item $u\cdot v=0$ implies that one of the vectors is zero.
\end{enumerate}
\begin{repere}{Self-assessment}I distinguish the norm from its square; I choose a suitable formula; I check the frame; I write an equivalence for a locus.

For each objective, indicate: achieved without help; achieved with help; to be reworked. Write one example and one mistake to avoid.\end{repere}
\endgroup
\ifprof\corrsection{3}\label{C03-CORR}
\textbf{Teacher use.} The answers below concern the tasks not worked in the lesson. First have students state their method; give graded help before the solution. The durations are estimates to be calibrated in class.
\subsection*{Diagnostic and remediation}
\begin{enumerate}
\item $(3;-4)$. \textit{Action:} Subtract end minus start on two drawings.
\item $\sqrt{3^2+(-4)^2}=5$. \textit{Action:} Review Pythagoras and distinguish norm from coordinates.
\item $\vect{AC}$. \textit{Action:} Draw the composite displacement.
\item $a^2+2ab+b^2$. \textit{Action:} Identify the cross term.
\item $\dfrac{3}{5}$. \textit{Action:} Name the sides relative to the angle.
\item $(1;2)$. \textit{Action:} Average the two coordinates separately.
\end{enumerate}
\subsection*{Investigation activity}
Take points on the circle with diameter AB; proof via the midpoint I: product $MI^2-9$. The endpoints are included because one of the vectors is zero.
\subsection*{Immediate applications}
\begin{enumerate}
\item $0$ then $-10$. The cosine is 0 and $-1$ respectively.
\item $3t-12=0$, so $t=4$.
\item The square is $4-2+9=11$; the norm is $\sqrt{11}$.
\item $32=16+16-32\cos A$, so $\cos A=0$ and $A=90^\circ$.
\item $MI^2-9=7$: circle with centre $(3;0)$ and radius 4.
\end{enumerate}
\subsection*{Exercises: method and results}
\paragraph{Exercise 3.1 -- Direct calculations.} Product $4-6=-2$; norms $\sqrt5$ and 5; sum $(5;1)$ with square 26. Check: $5+2(-2)+25=26$.
\paragraph{Exercise 3.2 -- A right triangle?} $\vect{AB}=(3;1)$, $\vect{AC}=(-1;3)$; product 0. $AB=AC=\sqrt{10}$, area $\dfrac{AB\times AC}{2}=5$ area units.
\paragraph{Exercise 3.3 -- Length without a frame.} $\vect{AC}=\vect{AB}+\vect{AD}$ and $\vect{BD}=\vect{AD}-\vect{AB}$. Product $\dfrac{4\times 3}{2}=6$; $AC^2=16+9+12=37$ and $BD^2=16+9-12=13$.
\paragraph{Exercise 3.4 -- Recovering an angle.} $\cos \theta =-\dfrac{6}{12}=-\dfrac{1}{2}$; in $[0;180^\circ]$, the angle is $120^\circ$. The domain ensures uniqueness.
\paragraph{Exercise 3.5 -- A variable locus.} $MI^2=k+\dfrac{AB^2}{4}$. If the right-hand side is negative: empty set; zero: the single point I; positive: circle whose radius is its square root. The discussion bears on a squared distance.
\paragraph{Exercise 3.6 -- Work done by a force.} $W=\dfrac{50\times 4\times 1}{2}=100$ J. At $120^\circ$, $W=-100$ J. The sign distinguishes help to the displacement from opposition; it does not describe a negative distance.
\paragraph{Exercise 3.7 -- Contrapositive.} The direct property ``orthogonal implies equal squares'' follows from the expansion. Its contrapositive is exactly the statement. We can also write $2\vec u\cdot\vec v\ne0$.
\paragraph{Exercise 3.8 -- Delayed review.} Al-Kashi: $\dfrac{9+16-25}{2}=0$. Converse of Pythagoras: triangle right-angled at A, so zero product. Name the hypotheses of each method.
\paragraph{CQFD, exercise 9 p.~250, questions a and b.} a: $1\times3+2\times(-1)=1$. b: $(-1)\times12+3\times7=9$. The scalar product is a number, not a vector.
\paragraph{CQFD, exercise 11 p.~250, questions 1, 2 and 3.} $\vect{AB}=(4;6)$, $\vect{AC}=(-3;4)$, $I(-1;-2)$ and $\vect{CI}=(5;-1)$. Hence $\vect{AB}\cdot\vect{AC}=12$, $\vect{AB}\cdot\vect{CI}=14$. Chasles gives $\vect{AB}\cdot\vect{AI}=12+14=26$. Two checks: $\vect{AI}=(2;3)$ gives $8+18=26$; or $\vect{AI}=\dfrac12\vect{AB}$ gives $\dfrac12(16+36)=26$.
\subsection*{MCQ: answers and reasons}
\textbf{MCQ 1.}
\begin{enumerate}
\item A. $6-4=2$.
\item C. Square root of 25.
\item D. Scalar criterion.
\item B. $\dfrac{2\times 3}{2}=3$.
\item C. Expand the scalar square.
\end{enumerate}
\textbf{MCQ 2.}
\begin{enumerate}
\item D. Negative cosine.
\item A. Coordinates by projection.
\item C. $MI^2=\dfrac{AB^2}{4}$.
\item B. Bilinearity: $2(-1)3$.
\item D. $BC^2=8$.
\end{enumerate}
\subsection*{True or false}
\begin{enumerate}
\item False: $u\cdot(-u)=-\norm u^2<0$ for $u\ne0$.
\item True: its product with every vector is 0.
\item False: $(1;0)$ and $(0;1)$ are nonzero and orthogonal.
\end{enumerate}
\paragraph{Differentiated teaching.} For a technical block, offer only the first step (domain, factor or reference sample space), then withdraw the help. For an independent student, ask for a second method and the treatment of a limiting case. Orally, allow an initial reformulation in the language the student has mastered, then require a precise mathematical sentence in French.
\fi
